\section{Model and statements}\label{sec:model}

\subsection{Cubical code and its two sectors}

Let \(K_L\) be the cubical cellulation of
\((\mathbb R/L\Z)^4\), with vertices \(\Z_L^4\) and unit coordinate
cells. Write \(C_j(K_L;\F)\) for its cellular \(j\)-chains and
\(\partial_j\) for the boundary. There is a physical qubit on each
plaquette, so the number of qubits is \(6L^4\). For each edge \(e\) and
cube \(c\), set
\begin{equation}
 A_e=\prod_{p\supset e}X_p,\qquad
 B_c=\prod_{p\subset\partial c}Z_p.
\end{equation}
These checks commute: an edge and a cube have either zero or two common
incident boundary plaquettes. Their joint \(+1\) eigenspace is the code
space. A \(Z\)-error is a chain \(E\in C_2(K_L;\F)\), with syndrome
\(s=\partial_2E\) on the edge checks. Since
\(\partial_1\partial_2=0\), every syndrome vertex has even degree.
A closed residual \(Z\)-chain is trivial on the code precisely when
it belongs to \(\operatorname{im}\partial_3\).

The dual cubical cellulation \(K_L^*\) identifies each primal
plaquette \(p\) with a dual plaquette \(p^*\), and each primal cube with
a dual edge. An \(X\)-error is consequently a dual 2-chain with a dual
edge syndrome. The two descriptions have the same geometry. The
dimension of \(H_2(K_L;\F)\) is \(\binom42=6\); equivalently the
coordinate 2-tori give six independent membrane classes. One may see
this from the product cellular complex of four circles: its mod-two
differentials vanish and its degree-two cells are indexed by pairs of
circle factors. Thus the code carries six logical qubits.

\subsection{The upper-corner rule}

On an integer lift, denote by \(e_i\) the coordinate unit vectors and
write
\[
 \operatorname{ht}(v)=v_1+v_2+v_3+v_4.
\]
For \(i<j\), the face with upper corner \(v\) has vertices
\[
 v,\quad v-e_i,\quad v-e_j,\quad v-e_i-e_j.
\]
Its tested edges are \([v-e_i,v]\) and \([v-e_j,v]\). A correction
attempt flips that plaquette exactly when both tested edges belong to
the syndrome. All six face types at every upper vertex have independent
rate-one Poisson attempt clocks. Ineffective rings remain part of the
graphical construction. On the dual lattice use the same upper-corner
convention in dual coordinates.

This is the convention in Drechsler--Oreg--Aharonov,
Definition 6.2 and Eq.~(169) \cite{DOA}. In the notation of
Pastawski--Clemente--Cirac, Eq.~(33), the tested sides leave the lower
corner \cite{PCC}. The reflection
\begin{equation}\label{eq:reflection}
 \mathfrak r:\Z_L^4\longrightarrow\Z_L^4,\qquad
             \mathfrak r(x)=-x
\end{equation}
takes our face with upper vertex \(v\) to their face with lower vertex
\(-v\), and takes its tested incoming edges to their outgoing lower
sides. The same applies to the dual lattice after choosing its
coordinate origin. Check Pauli names are conventional: their edge
checks are \(Z\)-type, whereas ours are \(X\)-type, so identifying the
two displayed stabilizer conventions also exchanges \(X\) and \(Z\)
uniformly, implementable by a Hadamard on every qubit. Reflection
preserves incidence, locality and the independent-clock laws. This
identifies the geometric rule up to reflection and the stated Pauli
convention.

\subsection{Time-independent generator and noise}

For a primal plaquette \(p\) with tested edges \(e,e'\), and its dual
plaquette with tested dual edges corresponding to cubes \(c,c'\), put
\begin{equation}
 \Pi_p^X=\frac{1-A_e}{2}\frac{1-A_{e'}}{2},\qquad
 \Pi_p^Z=\frac{1-B_c}{2}\frac{1-B_{c'}}{2}.
\end{equation}
Products are unambiguous because all checks commute. Define
\(\Diss[J](\rho)=J\rho J^\dagger-
\frac12(J^\dagger J\rho+\rho J^\dagger J)\). The model is
\begin{equation}\label{eq:generator}
 \Lind_L=
 \sum_p\bigl(\Diss[Z_p\Pi_p^X]+\Diss[X_p\Pi_p^Z]\bigr)
 +\sum_p\sum_{P\in\{X,Y,Z\}}\gamma_{p,P}\Diss[P_p],
 \qquad 0\le\gamma_{p,P}\le\eta .
\end{equation}
The coefficients are constant in time. Homogeneous \(\gamma_{p,P}\)
give translation invariance. The support of each jump has a bounded
diameter independent of \(L\). The Hamiltonian is zero.

The noise clocks for all pairs \((p,P)\), and the two families of
correction attempt clocks, are mutually independent. A noise ring
applies its single-qubit Pauli without a syndrome test. In the primal
\(Z\)-error marginal the noise rate is \(\gamma_{p,Z}+\gamma_{p,Y}\le
2\eta\), and in the dual \(X\)-error marginal it is
\(\gamma_{p,X}+\gamma_{p,Y}\le2\eta\). A \(Y\) mark belongs to both
marginals; their independence is neither asserted nor needed.

Correction rate one fixes the unit of physical time. If all correction
attempt rates instead equal \(\lambda>0\), rescaling time by
\(\lambda\) replaces \(\gamma_{p,P}\) by
\(\gamma_{p,P}/\lambda\). All statements below then apply with
\(\eta=\max_{p,P}\gamma_{p,P}/\lambda\) and horizon
\(\lambda^{-1}e^{cL}\).

\subsection{Readout and fidelity convention}

Choose a code encoding isometry \(V_L\) for the six logical qubits.
For the designated qubit \(Q\), prepare the other five logical qubits
in any fixed density matrix \(\sigma\). Write
\[
 \enc(\rho)=V_L(\rho\otimes\sigma)V_L^\dagger.
\]
The recovery \(\readout\) measures the current commuting checks,
applies fixed compatible primal and dual syndrome fillings, inverts
the encoding, and discards the other five logical qubits.
Section~\ref{sec:decoder} defines the fillings on every attainable
syndrome. Define
\begin{equation}\label{eq:fidelity}
 \logical=\readout\circ e^{t\Lind_L}\circ\enc,\qquad
 F_{\min}(L,t)=
 \min_{\|\psi\|=1}\langle\psi|
       \logical(|\psi\rangle\langle\psi|)|\psi\rangle.
\end{equation}
This is the squared-overlap convention of
Drechsler--Oreg--Aharonov, Definition 2.19. Their Definition 2.20 calls
the system an \(\epsilon\)-memory on an interval \(I\) when
\(\inf_{t\in I}F_{\min}(L,t)\ge1-\epsilon\) \cite{DOA}.
The recovery in this definition is fixed as a function of the measured
syndrome. It needs no earlier syndrome record.

\begin{theorem}[Finite-volume quantum memory]\label{thm:main}
There exist constants \(\eta_0,c,C>0\), an integer \(L_0\), and for
each integer \(L\ge L_0\) a syndrome-only recovery \(\readout\), such
that the following holds. For every set of constant rates
\(0\le\gamma_{p,P}\le\eta_0\), every choice of designated logical qubit
and fixed state \(\sigma\) of the other five qubits, and every
\(0\le t\le e^{cL}\), the dynamics \eqref{eq:generator}, initialized
by \(\enc\), satisfies
\begin{equation}\label{eq:main}
 \inf_{0\le t\le e^{cL}}F_{\min}(L,t)\ge1-Ce^{-cL}.
\end{equation}
The constants are uniform over these rates and inputs. In fact, for
every pure state \(|\Psi\rangle_{QR}\) with an arbitrary reference \(R\),
\[
 \langle\Psi|(\logical\otimes\id_R)
          (|\Psi\rangle\langle\Psi|)|\Psi\rangle
 \ge1-Ce^{-cL}
\]
through the same interval. The recovery can be chosen independently
of the rates, \(t\), \(\sigma\), and the input state.
\end{theorem}

The proof is completed in Section~\ref{sec:quantum}. A single event
will recover all six logical qubits, which explains the stated
uniformity in the unused logical state. The restriction to an encoded
start includes the reference-entangled version.

\begin{theorem}[Infinite-volume syndrome stability]\label{thm:infinite}
For either sector on \(\Z^4\), the independent-clock graphical process
started with empty syndrome is well defined. There exist
\(\eta_0,C,\mu>0\) such that, for every choice of constant
single-qubit Pauli rates bounded by \(\eta_0\), every vertex \(x\), every
deterministic \(t\ge0\), and every real \(d\ge0\),
\begin{equation}\label{eq:infinite-main}
 \Prb\{x\text{ belongs at time }t
   \text{ to a syndrome component of diameter }\ge d\}
       \le Ce^{-\mu d}.
\end{equation}
The same bound, after a fixed enlargement of \(C\), holds if the
component and its membership of \(x\) are observed at any common
instant of the slab of the proof containing \(t\).
Moreover, writing \(\Prb_\gamma\) for an allowed family of rates,
\[
 \lim_{\eta\downarrow0}
 \sup_{\substack{0\le\gamma_{p,P}\le\eta\\x\in\Z^4,\ t\ge0}}
 \Prb_\gamma\{x\text{ is occupied at time }t\}=0.
\]
Diameter is the ambient lattice sup diameter. The assertions are
marginal statements for a named sector.
\end{theorem}

Section~\ref{sec:infinite} proves this theorem, including existence of
the process. At \(t=0\) the syndrome is empty. In statements about a
whole slab we use the slab \([0,T]\) for \(t=0\), and otherwise
\(((k-1)T,kT]\) when \((k-1)T<t\le kT\).

Theorems~\ref{thm:main} and \ref{thm:infinite} have the independent
single-qubit Pauli scope stated above.
Theorem~\ref{thm:correlated} treats correlated stochastic Pauli faults.
The strongest statement, Theorem~\ref{thm:general}, covers arbitrary
sufficiently weak additive finite-range Markovian Lindblad noise,
including coherent and non-Pauli terms, with the exact correction.
Perturbations of the correction's own jump operators and general
non-Markovian baths remain outside that theorem.
The estimates yield existence of positive constants, rather than a
numerically useful threshold or a practical minimum size.
